\documentclass{article}
\usepackage{amsmath, amssymb}
\usepackage{a4wide}
\usepackage{enumitem}

\title{Stat 17 Week 2}
\author{Xiao Zhang}
\date{January 2026}

\begin{document}
\maketitle

\section*{Part I: Data, Variables, and Sampling}

\begin{itemize}

\item \textbf{Statistics and Data}
\begin{itemize}
  \item Statistics is the science of collecting, organizing, analyzing, and interpreting data
  \item Data consist of observations on individuals or objects
\end{itemize}

\item \textbf{Population, Sample, and Parameters}
\begin{itemize}
  \item \textbf{Population}: the entire group of interest
  \item \textbf{Sample}: a subset of the population actually observed
  \item A \textbf{parameter} describes a population (usually unknown)
  \item A \textbf{statistic} summarizes a sample and estimates a parameter
\end{itemize}

\item \textbf{Types of Variables}
\begin{itemize}
  \item \textbf{Categorical/qualitative} variables place individuals into categories
  \begin{itemize}
    \item Nominal: categories without order
    \item Ordinal: categories with a natural order
    \item Binary: two possible outcomes
  \end{itemize}
  \item \textbf{Quantitative} variables take numerical values
  \begin{itemize}
    \item Discrete: result of counting
    \item Continuous: result of measuring
  \end{itemize}
\end{itemize}

\item \textbf{Sampling Methods}
\begin{itemize}
  \item \textbf{Simple random sampling}: each individual has an equal chance of selection
  \item \textbf{Systematic sampling}: select every $k$th individual from a list
  \item \textbf{Stratified sampling}: divide population into groups, then randomly sample from each group
  \item \textbf{Cluster sampling}: randomly select groups and include all individuals within them
  \item \textbf{Convenience sampling}: easy to obtain but often biased
\end{itemize}

\item \textbf{Bias and Ethics}
\begin{itemize}
  \item Bias occurs when the sample is not representative of the population
  \item Common sources include sampling bias, response bias, and measurement bias
  \item Ethical data collection requires informed consent, privacy protection, and avoiding harm
\end{itemize}

\end{itemize}

\section*{Part II: Descriptive Statistics and Data Visualization}

\begin{itemize}

\item \textbf{Why We Summarize Data}
\begin{itemize}
  \item Large datasets are difficult to interpret without summaries
  \item Numerical and graphical summaries reveal patterns and unusual features
\end{itemize}

\item \textbf{Displaying Categorical Data}
\begin{itemize}
  \item Bar charts compare frequencies or relative frequencies
  \item Pie charts show proportions of a whole (best with few categories)
\end{itemize}

\item \textbf{Displaying Quantitative Data}
\begin{itemize}
  \item Histograms display the distribution of numerical data
  \item Useful for identifying shape, center, spread, and outliers
\end{itemize}

\item \textbf{Describing Distribution Shape}
\begin{itemize}
  \item Symmetric vs.\ skewed distributions
  \item Right-skewed: long tail to the right
  \item Left-skewed: long tail to the left
  \item Outliers, gaps, and clusters may appear
\end{itemize}

\item \textbf{Measures of Center}
\begin{itemize}
  \item Mean: arithmetic average
  \[
  \bar{x} = \frac{1}{n}\sum_{i=1}^n x_i
  \]
  \item Median: middle value of ordered data
  \item Mode: most frequent value (especially for categorical data)
\end{itemize}

\item \textbf{Mean vs.\ Median}
\begin{itemize}
  \item Symmetric distribution: mean $\approx$ median
  \item Right-skewed: mean $>$ median
  \item Left-skewed: mean $<$ median
  \item Median is more resistant to outliers than the mean
\end{itemize}

\item \textbf{Percentiles and Quartiles}
\begin{itemize}
  \item The $p$th percentile is the value below which $p\%$ of observations fall
  \item The median is the 50th percentile
  \item Quartiles divide ordered data into four equal parts:
  \begin{itemize}
    \item $Q_1$: 25th percentile
    \item $Q_2$: 50th percentile (median)
    \item $Q_3$: 75th percentile
  \end{itemize}
\end{itemize}

\item \textbf{Interquartile Range (IQR)}
\begin{itemize}
  \item The IQR measures the spread of the middle 50\% of the data
  \[
  \text{IQR} = Q_3 - Q_1
  \]
  \item The IQR is resistant to outliers
\end{itemize}

\item \textbf{Outliers and the 1.5$\times$IQR Rule}
\begin{itemize}
  \item Lower fence: $Q_1 - 1.5(\text{IQR})$
  \item Upper fence: $Q_3 + 1.5(\text{IQR})$
  \item Observations outside these fences are considered outliers
\end{itemize}

\item \textbf{Measures of Spread}
\begin{itemize}
  \item Range: maximum $-$ minimum
  \item Variance and standard deviation measure overall variability
  \[
  s^2 = \frac{1}{n-1}\sum_{i=1}^n (x_i-\bar{x})^2,
  \qquad
  s = \sqrt{s^2}
  \]
  \item Standard deviation has the same units as the original data
\end{itemize}

\item \textbf{The Complete Description of a Distribution}
\begin{itemize}
  \item A dataset should be described using:
  \begin{itemize}
    \item Shape
    \item Center
    \item Spread
  \end{itemize}
\end{itemize}

\end{itemize}

\end{document}
